\documentclass[12pt]{article}
\usepackage[utf8]{inputenc}
\usepackage[T1]{fontenc}
\usepackage{hyperref}
\title{Mang Language Resources for K-12 Curriculum Development}
\author{Compiled by Claude Sonnet 5.5 (claude-sonnet-5-5) for Andr\'as Kornai}
\date{}
\begin{document}
\maketitle

\section{Introduction}
Mang (Vietnamese: M\emph{a}ng; also Mang U; ISO 639-3 zng, Glottocode mang1378) is an Austroasiatic language spoken mainly in Lai Ch\^au Province, Vietnam, and across the border in Jinping County, Yunnan, China. It was first documented in 1974. Speaker counts in the sources range from about 2,200 to 4,900, and it is classed as severely endangered. Documentation is thin and no school materials were found. Mang is distinct from K\emph{h}ang (kjm), which some sources confuse with it.

\section{Primary Educational Resources}
No primer, textbook or school programme for Mang was found in search results.

\section{Linguistic and Reference Materials}
\begin{itemize}
\item Wikipedia article on the Mang language: \url{https://en.wikipedia.org/wiki/M%E1%BA%A3ng_language}
\item Wikipedia entry for ISO 639 zng: \url{https://wikipedia.com/wiki/ISO_639:zng}
\item Numeral system dataset for Mang (Yunnan, China, and Vietnam), Max Planck Institute: \url{https://lingweb.eva.mpg.de/channumerals/Mang.htm}
\item Verbix page for Mang verbs: \url{https://docs.verbix.com/Languages/Mang}
\end{itemize}

\section{Religious and Community Resources}
Joshua Project profile of the Mang in Vietnam: \url{https://joshuaproject.net/index.php/people_groups/13495/VM}

\section{Digital and Audio Resources}
\begin{itemize}
\item Endangered Languages Project sample entry for Mang: \url{https://www.endangeredlanguages.com/lang/1236/samples/10372}
\end{itemize}
The search summary also mentioned a Mang corpus in the Pangloss Collection (LACITO-CNRS), deposited by Michel Ferlus with speakers from P\'a Bon village, Lai Ch\^au, but no direct URL for it appeared in the results, so none is given.

\section{Curriculum Development Context}
With only a few thousand speakers, an endangered status, and no known school use, any curriculum would be a revitalization effort. It would need recording and documentation work first, then an agreed orthography, in close partnership with the community. Vietnamese (and Chinese in Yunnan) are the languages of schooling. The items above are reference sources, not teaching materials.

\end{document}